% Appended finite obstruction; does not assert an ultraproduct theorem.
\begin{theorem}[Proved: three-arrow obstructions in every cocycle degree]
\label{C4:three-arrow-high-degree}
For every integer $r\geq2$ there are finite objects $P_r,Q_r$ of
$\mathcal A_r(\mathbb Q)$ and three increasing embeddings
$f_0,f_1,f_2:P_r\to Q_r$ which have no rational affine equalizer.
Moreover, their finite simplicial coequalizer
\[
 R_r=\Delta^{Q_r}/(f_0\Delta^{P_r}=f_1\Delta^{P_r}
                         =f_2\Delta^{P_r})
\]
has nonzero rational homology in degree $r$. Thus there is no
bound, depending only on the number of arrows, on the homology
degree of such finite coequalizers or of their compatible-primitive
obstructions. The assertion uses three arrows independently of $r$;
the sizes of the two finite objects are allowed to grow with $r$.
\end{theorem}
\begin{proof}
We give the objects, cocycle, and detecting cycle explicitly. Let
\[
 A=\bigoplus_{i\in\mathbb Z}\mathbb Z e_i,\qquad
 \Gamma=A\rtimes\mathbb Z,
 \qquad(f,m)(g,n)=(f+T^m g,m+n),\quad Te_i=e_{i+1}.
\]
Write $a=(e_0,0)$ and $t=(0,1)$. These two elements generate
$\Gamma$. Give $\Gamma$ its left-invariant total order whose positive
cone consists of $(f,m)$ with $m>0$, or with $m=0$ and the highest
nonzero coordinate of $f$ positive. Indeed the indicated order of
$A$ is invariant under every shift, so this lexicographic extension
is a group order; conjugation also preserves its positive cone.
In particular every $a_i=(e_i,0)=t^iat^{-i}$ is positive. All finite
subsets below carry the induced order.

Define a shift-invariant alternating $r$-linear form $\omega$ on
$\mathbb Q\otimes A$ by setting
\[
 \omega(e_{i_1},\ldots,e_{i_r})=
 \begin{cases}
  \operatorname{sgn}(\pi),&
   (i_{\pi(1)},\ldots,i_{\pi(r)})=(j,j+1,\ldots,j+r-1)
   \text{ for some }j,\\
  0,&\text{otherwise}.
 \end{cases}
 \tag{TA.1}
\]
Here the sign is that of the sorting permutation, and a repeated
index gives zero. This specifies a well-defined form because every
input has finite support. On vertices $g_j=(v_j,m_j)\in\Gamma$ put
\[
 c(g_0,\ldots,g_r)
   =\omega(v_1-v_0,\ldots,v_r-v_0).
 \tag{TA.2}
\]
This is an alternating cochain on the full simplex with vertex set
$\Gamma$. Its coboundary is zero: expansion by multilinearity in
the alternating sum of its $r+2$ faces cancels every term. Equivalently,
the affine volume of the boundary of any $(r+1)$-simplex is zero.
Left multiplication by $(h,k)$ translates all $v_j$ by $h$ and
shifts their differences by $T^k$, so (TA.1) proves that $c$ is
$\Gamma$-invariant.

Let $V$ be the $2^r$ lamp vertices
\[
 V=\left\{\left(\sum_{i=0}^{r-1}\epsilon_i e_i,0\right):
                    \epsilon_i\in\{0,1\}\right\}.
\]
For each $0\leq i<r$, let $V_i^-$ be its face with $\epsilon_i=0$;
the opposite face is $V_i^+=a_i V_i^-$. Realize the word
$a_i=t^iat^{-i}$ as a sequence of left translations: first $i$
steps by $t^{-1}$, then one by $a$, then $i$ by $t$.
Let $P_0$ be the union of $V$ and all intermediate translates of
all the $V_i^-$ along these finitely many paths, including both
endpoints. Put
\[
 Q=P_0\cup aP_0\cup tP_0,\qquad
 P=\{g\in Q:ag\in Q\text{ and }tg\in Q\}.
 \tag{TA.3}
\]
Both are finite, $P_0\subset P\subset Q$, and $P$ is nonempty.
The three increasing injections $f_0(g)=g$, $f_1(g)=ag$,
$f_2(g)=tg$ have equal pullbacks of $c|_Q$, by invariance. Hence
$Q$ with cocycle $c|_Q$, and $P$ with their common pullback, are
objects of $\mathcal A_r(\mathbb Q)$ and these are morphisms.

In the oriented simplex chain group on $V$, use the standard cube
triangulation
\[
 z=\sum_{\pi\in S_r}\operatorname{sgn}(\pi)
 \left[0,e_{\pi(0)},e_{\pi(0)}+e_{\pi(1)},\ldots,
                         e_0+\cdots+e_{r-1}\right].
 \tag{TA.4}
\]
Lamp vectors in this formula mean their vertices in $\Gamma$.
Every displayed tuple is increasing, since its consecutive
differences are positive lamps. Its boundary is the sum of the
oriented triangulations of opposite cube faces, with opposite
signs. Each upper face is identified with its lower face in
$R_r$: along a step by $a$ or $t$ the starting face lies in $P_0$
and is identified with its translate by the corresponding
coequalizer relation; along a step by $t^{-1}$ the ending face
lies in $P_0$, and the same $t$-relation is used in reverse.
The paths specified above therefore identify the whole ordered
face simplices, not merely their vertices. Consequently the image
$\bar z$ of $z$ in the normalized simplicial chains of $R_r$ is a
cycle.

Because the restrictions of $c$ along all three arrows agree,
$c|_Q$ descends to an $r$-cocycle $\bar c$ on $R_r$. Formula
(TA.1) gives
\[
 c\left(0,e_{\pi(0)},e_{\pi(0)}+e_{\pi(1)},\ldots,
                  e_0+\cdots+e_{r-1}\right)
       =\operatorname{sgn}(\pi),
 \qquad \langle\bar c,\bar z\rangle=r!\ne0.
 \tag{TA.5}
\]
Thus $[\bar z]\ne0$ in $H_r(R_r;\mathbb Q)$. In particular the
cocycle on $Q$ has no primitive whose three restrictions to $P$
agree: such a primitive would descend to $R_r$ and make the
pairing in (TA.5) zero. Proposition~\ref{C4:rational-primitive-criterion}
now rules out a rational affine equalizer (and therefore an
integral one).
\end{proof}

\paragraph{Scope of the compactness approach.}
This theorem supplies a fixed number of arrows for rational
nonexactness in arbitrarily high cocycle degree. The ordinary
category ultraproduct nevertheless fails modular exactness:
Theorem~\ref{C4:ultraproduct-exclusion} gives a two-arrow
zero-cocycle family whose affine witnesses over every nonzero
ring require unbounded support. The three-arrow family just
constructed itself inherits this obstruction. Indeed $P_0$
contains all pure shifts $(0,m)$ for $-(r-1)\leq m\leq0$:
use the path for $i=r-1$ starting at the zero vertex of its lower
face. The set $Q$ also contains $(0,1)$. Restricting $f_0,f_2$
to these $r$ vertices gives the first-window and last-window
inclusions into the $r+1$ pure-shift vertices. Their cocycles
are zero. Any affine equalizer of the three arrows therefore
restricts to an affine equalizer of this window pair, with no
increase in support. Its support has size at least $r$, by
\eqref{C4:window-support}. Consequently even a selected subcategory of the
$r\to\infty$ ultraproduct which retains this distinguished
family has no affine equalizer for it over any nonzero ring.
That direct compactness construction cannot give the
total-category counterexample sought in C.4.
The present theorem does not contradict
Theorem~\ref{C4:degree-colimit-exclusion}, whose degree bound
depends on the size of a fixed source target object and whose
transition maps are actual functors.
